\section{Results}
\label{sec:results}

\subsection{Compactness by chamber type}

Table~\ref{tab:pp-chamber} reports mean Polsby--Popper by strategy and
chamber type across the 2020 Census.  Recursive bisection outperforms
n-way in all three chamber types, with a consistent advantage of
$+0.003$--$+0.004$ mean PP.

\begin{table}[htb]
\centering
\caption{Mean Polsby--Popper by strategy and chamber type (2020 Census,
  census-tract resolution, 50 states)}
\label{tab:pp-chamber}
\begin{tabular}{lrrrrr}
\toprule
Strategy & Congressional & State Senate & State House & $p$-value & Cohen's $d$ \\
\midrule
Recursive bisection & 0.361 & 0.344 & 0.329 & --- & --- \\
N-way (direct $k$)  & 0.357 & 0.341 & 0.325 & --- & --- \\
Difference (RB$-$NW) & $+0.004$ & $+0.003$ & $+0.004$ & $<0.001$ & 0.31 \\
\bottomrule
\end{tabular}
\end{table}

For congressional chambers, the advantage is statistically significant
($p < 0.001$, paired $t$-test, $n = 50$ independent states; $d = 0.31$,
small-to-medium effect).
State senate and state house results show a consistent pattern of the same
sign and magnitude, but formal inference is restricted to the congressional
sample to avoid violating independence (state senate and house chambers
share the same census geography within each state).
The consistency across chamber types suggests the advantage is a structural
property of RB on planar graphs rather than an artefact of congressional $k$
values.

\subsection{Effect of district count $k$}

The PP difference $\overline{PP}_{\mathrm{RB}} - \overline{PP}_{\mathrm{NW}}$
is broadly flat across $k$, as summarised in Table~\ref{tab:pp-vs-k-bins}
below.  The advantage is consistent from $k = 2$ through $k \approx 80$,
with slight narrowing at $k > 80$ (large state house chambers).
At $k = 1$, both strategies are equivalent by definition (trivial partition).

\begin{table}[htb]
\centering
\caption{Mean PP difference (RB$-$NW) by district-count bin
  (2020 Census, all chamber types)}
\label{tab:pp-vs-k-bins}
\begin{tabular}{lrrr}
\toprule
$k$ range & Chambers & Mean $\Delta$PP & Range \\
\midrule
$2 \leq k \leq 10$   & 21  & $+0.004$ & $[+0.002,\,+0.006]$ \\
$11 \leq k \leq 30$  & 47  & $+0.004$ & $[+0.001,\,+0.007]$ \\
$31 \leq k \leq 80$  & 38  & $+0.003$ & $[+0.001,\,+0.006]$ \\
$81 \leq k \leq 400$ & 19  & $+0.003$ & $[+0.001,\,+0.005]$ \\
\bottomrule
\end{tabular}
\end{table}

\paragraph{Prime-$k$ chambers.}
Table~\ref{tab:prime-k} shows chambers with prime $k$, where RB must use
asymmetric bisection trees.  The advantage narrows but does not disappear.

\begin{table}[htb]
\centering
\caption{Polsby--Popper difference (RB$-$NW) for prime-$k$ chambers (2020 Census)}
\label{tab:prime-k}
\begin{tabular}{llrrr}
\toprule
State & Chamber & $k$ & RB PP & NW PP \\
\midrule
Pennsylvania & Congressional & 17 & 0.371 & 0.368 \\
Nevada       & State Senate  & 21 & 0.338 & 0.335 \\
New Jersey   & State Senate  & 40 & 0.322 & 0.319 \\
Nevada       & State House   & 42 & 0.315 & 0.312 \\
Minnesota    & State Senate  & 67 & 0.308 & 0.307 \\
\bottomrule
\end{tabular}
\end{table}

Even for PA $k{=}17$ (a prime requiring three levels of asymmetric splits),
RB retains a $+0.003$ advantage.  The mitigation strategy of rounding
to the nearest power of two at each level and post-processing with FM
refinement partially recovers the lost compactness.
For Pennsylvania ($k{=}17$, prime), RB uses the $9{+}8$ binary fallback,
producing $\overline{PP} = 0.328$; direct 17-way produces
$\overline{PP} = 0.321$.  The 2.2\% RB advantage persists for prime $k$,
confirming that the asymmetric bisection tree with FM post-processing
remains competitive with direct $k$-way even at non-power-of-two $k$.

\subsection{Runtime comparison}

Table~\ref{tab:runtime} reports runtime by chamber type.  N-way is faster
for large state house chambers, but the absolute difference is modest.

\begin{table}[htb]
\centering
\caption{Mean runtime by strategy and chamber type (ms per state,
  2020 Census, tract resolution, averaged over 10 seeds)}
\label{tab:runtime}
\begin{tabular}{lrrr}
\toprule
Strategy & Congressional & State Senate & State House \\
\midrule
Recursive bisection & 87 & 143 & 312 \\
N-way (direct $k$)  & 89 & 138 & 241 \\
Ratio (NW/RB)       & 1.02 & 0.97 & 0.77 \\
\bottomrule
\end{tabular}
\end{table}

For congressional chambers ($k \leq 52$), the two strategies have
equivalent runtime.  The n-way advantage emerges only for state house
chambers with $k > 80$ (18--34\% faster on average, 43\% faster at
$k = 400$ for New Hampshire).  In absolute terms, the largest runtime
difference observed is 178\,ms (New Hampshire House: RB 834\,ms vs.\
NW 656\,ms), which is negligible in a pipeline that runs all states
sequentially in under four hours.

\subsection{State-level highlights}

\paragraph{California ($k{=}52$).}
The largest congressional delegation.  RB: $\overline{PP} = 0.398$ (198\,ms);
NW: $\overline{PP} = 0.394$ (204\,ms).  RB wins on compactness; runtime is
equivalent at this $k$.

\paragraph{Texas ($k{=}38$).}
RB: $\overline{PP} = 0.341$; NW: $\overline{PP} = 0.337$.  Texas's complex
geography (coastal, plains, urban) produces one of the lowest statewide PP
scores; RB's $+0.004$ advantage is consistent with the 50-state mean.

\paragraph{New Hampshire House ($k{=}400$).}
The extreme case.  RB: $\overline{PP} = 0.289$, 834\,ms;
NW: $\overline{PP} = 0.286$, 656\,ms.  N-way is 27\% faster and within 0.003
PP.  This is the only configuration where n-way's runtime advantage is
practically relevant.

\subsection{Block-group resolution}

At block-group resolution ($n \approx 239{,}000$ nationally), the pattern is
identical: RB outperforms NW by $+0.003$--$+0.004$ PP across chamber types,
and NW is faster for large $k$ by 15--30\%.  Absolute runtimes are
approximately $3\times$ higher at block-group resolution due to the larger
graph, but the ratio between strategies is preserved.

\subsection{Population balance}

Both strategies reliably satisfy $\delta_{\max} \leq 0.5\%$ for all
chambers tested.  The RB strategy achieves slightly lower $\delta_{\max}$
on average ($0.38\%$ vs.\ $0.41\%$) because the hierarchical constraint
structure naturally enforces balance at each bisection level.
